\subsection{The remaining pure-band sources through order ten}

\begin{lemma}[Proved: row multiplicity in fourth homology]
\label{CP2:A-pure-band-fourth-row-multiplicity}
Let $M$ be the identity-adjoined $r\times c$ rectangular
band, with $r,c\ge2$, and let $N=R_3$.
For every finite nonempty left-$N$, right-$M$ biset $S$,
the rational fourth Betti number of its collage is
divisible by $r-1$.  In particular $r\ge3$ excludes
the rational homology of $\mathbb{CP}^2$, without
any state-count or boundedness hypothesis.
\end{lemma}
\begin{proof}
Put $R=\mathbb QM$, $T=\mathbb QN$, and $L=\mathbb Q[S]$.
Every core orbit in $SK$ is a quotient of the $c$
column labels, exactly as in
Lemma~\ref{CP2:A-two-by-three-core-partitions}.
Write $W_{\pi_O}$ for the kernel of the column-difference
map $\mathbb Q^{c-1}\to\mathbb Q[O]$; it has dimension
$c-|O|$.

Let $e=e_{00}$ and $q=1-e$.  The complete resolution
in Lemma~\ref{CP2:A-band-spheres} is
\[
 0\longrightarrow (Re)^{(r-1)(c-1)}
 \longrightarrow (Rq)^{r-1}\longrightarrow Re
 \longrightarrow\mathbb Q\longrightarrow0.
\]
The first differential is the direct sum, indexed by
the $r-1$ non-base rows, of the same column-difference
map $(Re)^{c-1}\to Rq$.
After tensoring with $L$, its kernel is therefore
\[
 \operatorname{Tor}^{R}_2(L,\mathbb Q)
   \cong\left(\bigoplus_{O\subseteq SK}W_{\pi_O}\right)^{r-1}
                                                               \tag{1}
\]
as a left-$T$ module.  This is a natural direct sum:
the commuting $N$ action acts on core orbits and column
vectors, and leaves the row-indexed multiplicity unchanged.
The right-$T$ projective-dimension-one bound and the
two-block sequence give, exactly as before,
\[
 H_4(B\mathcal C_S;\mathbb Q)
 \cong
 \operatorname{Tor}^{T}_1
       (\mathbb Q,\operatorname{Tor}^{R}_2(L,\mathbb Q)).
\]
Applying this additive functor to (1) proves the divisibility.
\end{proof}

\begin{theorem}[Proved: six states for the two-by-four band]
\label{CP2:A-two-by-four-band-six-states}
Let $M$ be the identity-adjoined $2\times4$ rectangular
band and let $N=R_3$.  If $S$ has at most six states,
its collage cannot have the rational cohomology ring
of $\mathbb{CP}^2$.
More precisely, if its fourth rational Betti number is
one, all its rational degree-two products vanish.
\end{theorem}
\begin{proof}
A singleton core orbit supplies an actual connecting
arrow fixed by $M$, and the fixed-arrow corollary
kills all products.  Assume that every core orbit
has at least two states.
Let $\mathcal O$ denote the finite set of core orbits.
The two actors of $N$ induce retractions on $\mathcal O$
with the same image, and an orbit image has a coarser
column partition.  A self-map of one core orbit that
preserves column labels is its identity.

There are at most three core orbits.
If there is only one, it is pointwise fixed by $N$,
so its column kernel is a trivial left-$N$ module
and its first $N$-Tor group is zero.
If there are three, all three have exactly two states.
An $N$ map between them must preserve the two-block
column partition.  Thus the module
$\bigoplus W_{\pi_O}$ is a direct sum, over the
occurring partitions, of two identical copies of
the permutation module on the corresponding orbit set.
Its first $N$-Tor dimension, and hence the fourth
Betti number of the collage, is even.
Neither case allows fourth Betti number one.

Suppose there are exactly two core orbits.
If both are fixed by $N$, the same Tor group is zero.
Otherwise both $N$ retractions map the unique nonfixed
orbit $O_1$ to the unique fixed orbit $O_0$.
Their maps on column kernels coincide with the inclusion
$W_{\pi_1}\hookrightarrow W_{\pi_0}$.
The explicit right-zero resolution identifies
\[
 \operatorname{Tor}^{T}_1
     (\mathbb Q,W_{\pi_0}\oplus W_{\pi_1})
      \cong W_{\pi_1}.
                                                               \tag{2}
\]
Indeed the kernel of either retraction on this direct
sum consists of the vectors $(-w,w)$ with
$w\in W_{\pi_1}$, and the difference of the two
retractions is zero.
Fourth Betti number one therefore forces
$\dim W_{\pi_1}=4-|O_1|=1$, so $|O_1|=3$.
The fixed orbit $O_0$ has two or three states and a
coarser partition.  In particular, the unique pair
of columns identified by $\pi_1$ is also identified
by $\pi_0$.

The actual subcollage on the connecting subset
$S_0=SK=O_0\sqcup O_1$ has the same second source-Tor
module as the original collage: no outside state
contributes to the kernel of the column-difference
map.  Its inclusion therefore induces an isomorphism
on rational fourth homology.  It suffices to kill
all degree-two products on this subcollage.

Interchange the identified pair of columns, fixing
the other two columns and the row labels.
Choose a fixed column as column zero in the projective
resolution.  This column transposition $\tau$ extends
to an actual involution of the subcollage, fixing $N$
and all its connecting states.
The resolution lift acts on the source-Tor group
by its ordinary permutation of column-difference
coordinates.  Its action on the line $W_{\pi_1}$
in (2) is $-1$.  Hence it acts as $-1$ on the
fourth homology of the subcollage.

For a union of core orbits, the right action is row
independent and the map
$(Le)^{c-1}\to Lq$ is surjective: the column differences
span all non-base states in every orbit.
Thus $\operatorname{Tor}^{R}_1(L,\mathbb Q)=0$,
and $\operatorname{Tor}^{R}_0(L,\mathbb Q)$ is the
permutation module on the core orbit set.
The column transposition acts as the identity on this
module, so the double-bar spectral sequence identifies
$H_1(E_{S_0};\mathbb Q)$ with
$\operatorname{Tor}^{T}_1(\mathbb Q,\mathbb Q[\mathcal O])$
and makes it entirely positive under $\tau$.

The negative eigenspace of $H_2(BM;\mathbb Q)$ is
one-dimensional, generated by the difference of the
two interchanged columns.  It lies in $W_{\pi_0}$.
Because $O_0$ is pointwise fixed by $N$, its actual
subbridge maps onto that negative line in $H_2(BM)$.
The negative part of the exact cohomology sequence
\[
 H^1(E_{S_0};\mathbb Q)\longrightarrow
 H^2(B\mathcal C_{S_0};\mathbb Q)\longrightarrow
 H^2(BM;\mathbb Q)\longrightarrow H^2(E_{S_0};\mathbb Q)
\]
therefore shows that all of $H^2(B\mathcal C_{S_0};\mathbb Q)$
is positive.  Every product of two second-degree
classes is positive, while fourth cohomology is negative;
all such products vanish.  The fourth-homology
isomorphism transfers this conclusion to the original
collage, proving the theorem.
\end{proof}

\begin{corollary}[Proved: the remaining pure-band catalog rows]
\label{CP2:A-pure-band-nine-ten-right-zero}
The identity-adjoined $2\times4$ band at positive
order-nine catalog index $8042$ is excluded from
the eighteen-arrow budget $(9,3,6)$.
The identity-adjoined $3\times3$ band at positive
order-ten index $157130$ is excluded from $(10,3,5)$.
\end{corollary}
\begin{proof}
Apply the six-state theorem to the first source and
the even fourth-Betti-number obstruction to the second.
Both literal catalog tables are among the six
independently reconstructed pure-band bindings in
\path{acyclic/outputs/refined_small_corner_sources_independent_review.json}.
\end{proof}
